\bookchapter{Приложение A: Пошаговое руководство по распаковке нового телефона}

\booksection{«Три шага» распаковки и проверки устройства:}

\begin{itemize}
  \item \textbf{Осмотрите внешний вид}: проверьте корпус и экран на наличие царапин и вмятин.
  \item \textbf{Сверьте информацию}: убедитесь, что IMEI/серийный номер на коробке совпадает с тем, что указан в чеке.
  \item \textbf{Проверьте функции}: сделайте звонок для проверки динамика и микрофона, сделайте фото для проверки камеры.
\end{itemize}

\booksection{«Ключевые напоминания» при активации и первоначальной настройке:}

\begin{itemize}
  \item \textbf{Сеть}: сначала подключитесь к стабильному Wi-Fi.
  \item \textbf{Аккаунт}: войдите в свой Apple ID, аккаунт Huawei, аккаунт Xiaomi и т. д.
  \item \textbf{Базовые настройки}: включите «Найти мой телефон» и установите надёжный пароль блокировки экрана.
  \item \textbf{Настройка eSIM}: если ваш телефон поддерживает eSIM (например, iPhone и некоторые Android-флагманы), можно добавить тариф eSIM в настройках без вставки физической SIM-карты. При поездке за границу eSIM особенно удобна: можно напрямую купить тариф eSIM у местного оператора.
  \item \textbf{Перенос данных}: при переносе данных со старого телефона убедитесь, что оба устройства достаточно заряжены. Для iPhone используйте «Быстрое начало», для Android — официальное приложение для переключения.
\end{itemize}

\booksection{«Безопасность прежде всего» при установке приложений:}

\begin{itemize}
  \item \textbf{Сначала официальный магазин}: для Apple используйте App Store, для Android — встроенный магазин приложений телефона.
  \item \textbf{Не раздавайте разрешения направо и налево}: когда при установке появляются запросы на разрешения, подумайте: «Действительно ли этому приложению это нужно?»
\end{itemize}
